\section{Introduction}
\label{sec:intro}

Flow matching~\cite{lipman2022flow, rectified_flow} models have become dominant in image generation~\cite{sd3, flux2024} due to their solid theoretical foundations and strong performance in producing high quality images. However, they often struggle with composing complex scenes involving multiple objects, attributes, and relationships~\cite{T2i-compbench, Gpt-imgeval}, as well as text rendering~\cite{textdiffuser}. At the same time, online reinforcement learning (RL)~\cite{sutton1998reinforcement} has proven highly effective in enhancing the reasoning capabilities of large language models (LLMs)~\cite{deepseek_r1, openai_o1}. While previous research has mainly focused on applying RL to early diffusion-based generative models~\cite{ddpo} and offline RL techniques like direct preference optimization~\cite{dpo} for flow-based generative models~\cite{videoalign, skyreels_v2}, the potential of online RL in advancing flow matching generative models remains largely unexplored. In this study, we explore how online RL can be leveraged to effectively improve flow matching models.


Training flow models with RL presents several critical challenges: (1) Flow models rely on a deterministic generative process based on ODEs~\cite{rectified_flow}, meaning they cannot sample stochastically during inference. In contrast, RL relies on stochastic sampling to explore the environment, learning by trying different actions and improving based on rewards. \textit{This need for stochasticity in RL conflicts with the deterministic nature of flow matching models.} 
(2) Online RL depends on efficient sampling to collect training data, but flow models typically require many iterative steps to generate each sample, limiting efficiency. This issue is more pronounced with large models~\cite{flux2024, sd3}. To make RL practical for tasks like image or video generation, \textit{improving sampling efficiency is essential}.
% (2) {Online RL requires efficient sampling to gather training data, but flow models typically involve many iterative steps to generate each sample, significantly reducing sampling efficiency.} This issue becomes even more critical with large, advanced models~\cite{flux2024, sd3}, which are computationally intensive. Therefore, to make RL effective in tasks such as image or video generation, \textit{improving sampling efficiency within RL training is crucial}.

% (3) RL excels in improving LLM reasoning ability because mathematical tasks only require final-answer correctness. In contrast, image generation lacks a singular quality metric. For instance, in object counting tasks, optimizing solely for accurate object counts using the correctness of generated object counts as the reward model frequently results in correct quantities but reduced overall image quality.


% To address these challenges, 
% in this work, we propose the \textbf{Flow-GRPO} approach to introduce the GPRO~\cite{grpo} into the flow matching models for alignment on text-to-image tasks,
% where two key strategies are introduced.
% First, to overcome the deterministic limitations of the original flow model, we employ the \textbf{ODE-to-SDE} strategy by converting the ODE-based flow model into a Stochastic Differential Equation (SDE) framework. Specifically, in the ODE-to-SDE strategy, we convert the deterministic Flow-ODE into an equivalent Flow-SDE that matches the original model's marginal distributions, introducing randomness into the model's generation process.
% Second, to achieve high sampling efficiency for online RL, we introduce the \textbf{Denoise Reduction} strategy by reducing the denoising steps for training data generation while retaining the original denoising steps during inference.
% The key insight is that significant fewer denoising steps during RL training may be enough. In our experiemnts, we found training with fewer denoising steps matches the performance obtained with the full schedule while markedly cutting data‑generation cost.

To address these challenges, we propose \textbf{Flow-GRPO}, which integrates GRPO~\cite{grpo} into flow matching models for text-to-image (T2I) generation, using two key strategies.
First, we adopt the \textbf{ODE-to-SDE} strategy to overcome the deterministic nature of the original flow model. By converting the ODE-based flow into an equivalent Stochastic Differential Equation (SDE) framework, we introduce randomness while preserving the original marginal distributions.
Second, to improve sampling efficiency in online RL, we apply the \textbf{Denoising Reduction} strategy, which reduces denoising steps during training while keeping the full schedule during inference. Our experiments show that using fewer steps maintains performance while significantly reducing data generation costs.

% This reduction in computational complexity significantly accelerates the training process without compromising the model's performance. 
% Our experiments
 % Our method, in contrast, 
% , which aims to  stochastic sampling 
% (1) SDE sampling: By converting the ODE-based flow model into a Stochastic Differential Equation (SDE) framework, we enable stochastic sampling. This conversion introduces randomness into the model's generation process, overcoming the deterministic limitations of the original flow model. This approach is distinct from the traditional strategy in diffusion tasks, where the SDE is typically converted into an ODE to achieve higher generation quality.
% (2) FastDenoise: We utilize significantly fewer denoising steps during both sample generation and training, while retaining the original denoising steps during inference. The key insight behind this reduction strategy is that training on fewer denoising steps can xxx. This reduction in computational complexity significantly accelerates the training process without compromising the model's performance. Our experiments verify this: fewer denoising steps during training yield results comparable to those achieved with larger denoising steps during training.
% (3) applying a KL constraint instead of a general human-preference reward model to effectively prevent noticeable declines in image quality.

% We empirically evaluate our Flow-GRPO on text-to-image tasks with various reward types. (1) Verifiable rewards, primarily using the GenEval~\cite{geneval} benchmark and visual text rendering task. For example, GenEval comprises compositional image generation tasks (e.g., generating specified object counts, colors, and spatial relationships), which can be automatically assessed via traditional detection methods. On GenEval, Flow-GRPO improves Stable Diffusion 3.5 Medium (SD3-M)~\cite{sd3}’s accuracy from 63\% to 95\% (Figure~\ref{fig:teaser}), surpassing the performance of the state-of-the-art GPT-4o~\cite{gpt4o} model. For the visual text rendering task, we improve SD3-M's accuracy from 56\% to 92\%, significantly enhancing the model’s ability to generate text. (2) Model-based rewards, such as human preference Pickscore~\cite{kirstain2023pick} reward, further demonstrate that our framework is task-independent, highlighting its generalizability and robustness. 
We evaluate Flow-GRPO on T2I tasks with various reward types. (1) Verifiable rewards, using the GenEval~\cite{geneval} benchmark and visual text rendering task. GenEval includes compositional image generation tasks (e.g., generating specific object counts, colors, and spatial relationships), which can be automatically assessed with object detection methods. Flow-GRPO improves the accuracy of Stable Diffusion 3.5 Medium (SD3.5-M)~\cite{sd3} from 63\% to 95\% on GenEval, outperforming the state-of-the-art GPT-4o~\cite{gpt4o} model. For visual text rendering, SD3.5-M’s accuracy increases from 59\% to 92\%, greatly enhancing its text generation ability. (2) Model-based rewards, such as the human preference Pickscore~\cite{kirstain2023pick} reward. These results show that our framework is task independent, demonstrating its generalizability and robustness. Importantly, all improvements are achieved with very little reward hacking, as demonstrated in Figure~\ref{fig:teaser}.

% To summarize, our contribution is summarized as follows:
% \begin{itemize} [leftmargin=15pt]
%   \item We are the first to introduce GRPO to flow matching models by converting ODE sampling into SDE sampling, demonstrating that online RL is highly effective for T2I tasks. Specifically, our Flow-GRPO improves SD3-M accuracy from 63\% to 95\% without compromising image quality.
%   \item We find that for online RL, it is not necessary to use the standard longer timesteps when collecting training samples. We can utilize significantly fewer denoising steps during sample generation, while retaining the original denoising steps during testing, which can significantly accelerate the training process.
%   \item We demonstrate that the Kullback-Leibler (KL)  constraint is highly effective in preventing reward hacking\footnote{The reward becomes large, but image quality significantly decreases} without sacrificing image quality in image generation. 
%   We also note that KL is not simply analogous to early stopping. By adding an appropriate KL, we can achieve the same high reward as the KL-free version, , just with longer training time.
% \end{itemize}

To summarize, the contributions of Flow-GRPO are as follows:
\begin{itemize}[leftmargin=15pt]
\item We are the first to introduce GRPO to flow matching models by converting deterministic ODE sampling into SDE sampling, showing the effectiveness of online RL for T2I tasks. Flow-GRPO improves SD3.5-M accuracy from 63\% to 95\% without noticeably compromising image quality.
\item We find that online RL for flow matching models does not require the standard long timesteps for training sample collection. By using fewer denoising steps during training and retaining the original steps during testing, we can significantly accelerate the training process.
% \item We find that online RL with flow matching models can generate training samples with fewer denoising steps, greatly accelerating training. Full denoising is used at test time to maintain performance.
\item We show that the Kullback-Leibler (KL) constraint effectively prevents reward hacking, where reward increases at the cost of image quality or diversity. KL regularization is not empirically equivalent to early stopping. With a proper KL term, we can match the high reward of the KL-free version while preserving image quality, albeit with longer training.
\end{itemize}
